%!TEX root = ../../../main.tex
\chapter{Cloud-Zentrierung und infrastrukturelle Abhängigkeiten}\label{chap:cloud-ursachen}
% Kennzeichnung der Autorenschaft (Vorgabe der Arbeit) -- Form ggf. in der Gruppe vereinheitlichen
\noindent\textit{Verfasser: [Name]}


Beim Smart-TV steht die Verwertung von Nutzungsdaten im Vordergrund. Die
Bindung an Herstellerserver betrifft jedoch auch die Frage, ob ein Gerät
überhaupt funktioniert. Dieses und das folgende Kapitel untersuchen daher,
warum Hersteller Smart-Home-Geräte an ihre Server binden und welche Folgen
das für Funktion, Lebensdauer und Kontrolle der Käufer hat. Die leitende
These lautet: Cloud-Abhängigkeit ist weniger eine technische Notwendigkeit
als eine Folge von Geschäftsmodellen.

\section{Cloud-zentrierte Architektur am Beispiel des Saugroboters}\label{sec:cloud-architektur}

Viele Smart-Home-Geräte führen zentrale Funktionen nicht selbst aus, sondern
über Benutzerkonto, Hersteller-App und Hersteller-Server. Hoofnagle, Kesari
und Perzanowski bezeichnen diese dauerhafte Verbindung zwischen verkauftem
Produkt und Anbieter als \enquote{Tethered Economy}: Der Hersteller kann
Funktionen auch nach dem Kauf verändern oder entziehen
\autocite[vgl.][]{hoofnagle:2019}. Nach Tusikov macht gerade die
Angewiesenheit auf Cloud-Dienste vernetzte Produkte für ihre Hersteller
besonders steuerbar, bis hin zum gezielten Unbrauchbarmachen, dem
sogenannten \enquote{Bricking} \autocite[vgl.][S.~2, 7]{tusikov:2019}.

Der Saugroboter veranschaulicht diese Architektur. Er erfasst die Wohnung
mit LiDAR- oder Kamerasensoren und erzeugt daraus Raumkarten. Kartenverwaltung,
raumweise Reinigung, Sperrzonen und Zeitpläne laufen bei vielen Modellen
jedoch über App und Cloud des Herstellers. Das Open-Source-Projekt Valetudo,
das in Kapitel~7 behandelt wird, versteht sich ausdrücklich als Ersatz für
diese Cloud und begründet dies mit Datenschutz, der Unabhängigkeit von
Internetausfällen und dem Schutz vor nachträglich eingeführten Abonnements
\autocite[vgl.][]{valetudo:2025}.

Daraus ergeben sich zwei Abhängigkeitsebenen. Die erste ist der Fortbestand
des Herstellers. Der Roomba-Hersteller iRobot beantragte im Dezember~2025
Gläubigerschutz nach US-Insolvenzrecht und ging vollständig an seinen
Auftragsfertiger Picea Robotics über \autocite[vgl.][]{fortune:2025}. App
und Support liefen zwar weiter \autocite[vgl.][]{trustedreviews:2026}, doch
die Kontrolle über Wohnungskarten und Dienste wechselte den Eigentümer, ohne
dass die Nutzer Einfluss darauf hatten. Die zweite Ebene ist die
Infrastruktur: Als Amazon Web Services im Oktober~2025 ausfiel, ließen sich
die vernetzten Betten des Herstellers Eight Sleep nicht mehr steuern. Dass
der Hersteller kurz darauf eine Steuerung ohne Cloud nachrüstete, zeigt, dass
die vollständige Abhängigkeit eine Designentscheidung war
\autocite[vgl.][]{patterson:2025}.

Die Cloud ist damit nicht pauschal abzulehnen; sie ermöglicht Fernzugriff,
einfache Einrichtung und zentrale Updates. Versteht man Souveränität im Sinne
von Kapitel~2 als informierte Selbstbestimmung und nicht als Autarkie
\autocite[vgl.][]{ternes:2022}, liegt das Problem im Fehlen einer lokalen
Rückfallebene.

\section{Vendor-Lock-in über Schnittstellen}\label{sec:cloud-lockin}

Sind Kunden auf ein System festgelegt, wird der Wechsel zu einem
Konkurrenzprodukt teuer, und der Anbieter gewinnt Spielraum bei Preisen und
Konditionen \autocite[vgl.][]{shapiro:1999}. Im Cloud-Kontext entsteht ein
solcher Lock-in vor allem durch proprietäre Schnittstellen und mangelnde
Interoperabilität \autocite[vgl.][]{oparamartins:2016}. Eine Studie im
Auftrag des Umweltbundesamts (\acrshort{UBA}) beschreibt das Bemühen der
Hersteller, Kunden über proprietäre Systeme an ihre Plattformen zu binden,
als wirtschaftlichen Zielkonflikt, der langlebigen Produkten entgegensteht
\autocite[vgl.][S.~39]{jaegererben:2023}.

Wie gezielt Hersteller diesen Hebel einsetzen, zeigt der Garagentorhersteller
Chamberlain. Er sperrte 2023 den Zugriff von Drittanbieter-Software auf sein
myQ-System, woraufhin die Open-Source-Plattform Home Assistant ihre seit 2017
bestehende Integration entfernen musste. Offizielle Partner zahlen
Chamberlain dagegen für den Zugang \autocite[vgl.][]{homeassistant:2023},
und die herstellereigene App zeigt Werbung, die Nutzern von Drittanbieter-Apps
entgeht \autocite[vgl.][]{ninetofivemac:2023}. Samsung verlangt seit
Oktober~2026 von privaten Entwicklern rund fünf US-Dollar monatlich für den
bislang freien Zugang zur Programmierschnittstelle (\acrshort{API}) von
SmartThings, während die eigene App kostenfrei bleibt
\autocites[vgl.][]{sammobile:2026}{androidpure:2026}. In beiden Fällen wird
keine neue Funktion bepreist oder geschützt, sondern die Möglichkeit, das
eigene Gerät außerhalb des Hersteller-Ökosystems zu steuern. Gerade diese
Interoperabilität würde einen Wechsel erleichtern.

\section{Hardware as a Service: nachträgliche Monetarisierung}\label{sec:cloud-haas}

Hinter dem Lock-in steht ein strukturelles Finanzierungsproblem: Das Gerät
wird einmal verkauft, der Serverbetrieb verursacht dauerhaft Kosten. Der
Schweizer Anbieter myStrom machte dies 2024 offen, als er bislang kostenlose
Funktionen in ein Abonnement überführte. Ohne diesen Schritt wäre das
Unternehmen nach eigener Aussage Anfang~2025 in Konkurs gegangen; die Marge
von rund fünf Franken pro Gerät decke die Kosten nur für etwa 18~Monate
\autocite[vgl.][]{itreseller:2024}. Auch der Thermostathersteller Tado
testete 2025 eine Gebühr für seine App \autocite[vgl.][]{heise:2025b}. Das
Gerät wird so zur Plattform für wiederkehrende Erlöse, zu denen auch die in
Kapitel~3 beschriebene Verwertung von Nutzungsdaten zählt.

Abonnements sind nicht per se illegitim, denn Serverkosten sind real. Ein
beim Kauf transparent ausgewiesenes Abo erlaubt eine informierte
Entscheidung. Problematisch ist die nachträgliche Einschränkung oder
Bepreisung von Funktionen bereits bezahlter Geräte: Dann bestimmt der
Hersteller und nicht der Käufer den Funktionsumfang, was die
Entscheidungssouveränität unmittelbar berührt.

Dass der Markt dies nicht korrigiert, liegt an einer Informationsasymmetrie.
Die US-amerikanische Federal Trade Commission (\acrshort{FTC}) fand 2024 auf
den Herstellerseiten von 161 der 184 untersuchten vernetzten Produkte keine
Angabe zur Dauer der Software-Unterstützung \autocite[vgl.][]{ftc:2024}.
Können Käufer die Qualität eines Gutes nicht beurteilen, entscheiden sie nach
dem Preis, und hochwertige Anbieter werden nicht belohnt
\autocite[vgl.][]{akerlof:1970}. Die Lebensdauer der Herstellerdienste ist ein
solches verdecktes Qualitätsmerkmal, weshalb Transparenzpflichten wie das in
Kapitel~9 behandelte IT-Sicherheitskennzeichen ansetzen.
